\subsection{Simplifying Expressions: Basic Moves} 
Recall:
\begin{definition}
\begin{itemize}
\item An \term{expression} is a sequence of numbers, variables, and symbols that can be resolved to represent a number.
\item To \term{evaluate} an expression is to determine what number it represents.
\item Two expressions are said to be \term{equivalent} if they represent the same number, no matter what any variables represent. To \term{rewrite} an expression is to write an equivalent expression in a different form.
\item To \term{simplify} an expression is to write an equivalent expression in some standard ``simplest'' from.
\item The \term{terms} of an expression are parts of an expression combined using \emph{addition}. The \term{factors} of an expression are parts of an expression combined using \emph{multiplication}.
\end{itemize}
\end{definition}
There are a number of common steps we use to simplify expressions.
\begin{itemize}
\item Cancel common factors in the numerator and denominator of a fraction.
\item Use the commutative property to rearrange terms and factors.
\item Multiply numbers within terms.
\item Use the distributive property to eliminate parentheses around addition\\(or \makeblank).
\item ``Combine Like Terms'' (using the distributive property).
\item Use properties of exponents to combine factors.
\end{itemize}
\begin{Example}{Simplify $3(x+4)-5(x-3)$.}
\why{5}
\end{Example}